\clearpage\pagestyle{otMain}\setlength{\otBlocEspacement}{2.5mm}\fontsize{11}{15}\selectfont\corrigesChapitre{1}{Entiers et décimaux}
\begin{blocLivre}{cor-1-1}\begin{corrige}{1}{Lire un grand nombre}\demarche{Lire et écrire classe par classe.}
On lit classe par classe : « trois-milliards-cinquante-deux-millions-quatre-cent-mille ». Le second nombre s’écrit $\res{12\,500\,000\,000}$ : 12 dans la classe des milliards, 500 dans celle des millions.
\end{corrige}\end{blocLivre}
\begin{blocLivre}{cor-1-2}\begin{corrige}{2}{Chiffre ou nombre}\demarche{Pour un nombre de…, garder tout ce qui est à gauche du rang.}
Chiffre des millions : $\res2$ ; nombre de millions : $\res{7\,342}$. Chiffre des centaines de mille : $\res0$ ; nombre de centaines de mille : $\res{73\,420}$.
\end{corrige}\end{blocLivre}
\begin{blocLivre}{cor-1-3}\begin{corrige}{3}{Unités de numération}\demarche{Chaque rang vaut dix fois le rang de droite.}
$1=\res{10}$ dixièmes ; 1 dixième $=\res{10}$ centièmes ; 1 centième $=\res{10}$ millièmes ; $1=\res{1\,000}$ millièmes ; 1 millier $=\res{10}$ centaines.
\end{corrige}\end{blocLivre}
\begin{blocLivre}{cor-1-4}\begin{corrige}{4}{Lire une grille}\demarche{Compter les unités, puis les dixièmes, puis les centièmes.}
Deux unités entières, 4 colonnes de dixièmes et 7 petits carrés de centièmes : $2+\frac4{10}+\frac7{100}=\res{\frac{247}{100}}=\res{2{,}47}$.
\label{fig-cor-1-4}\figureGeo[corrige]{num-4}
\end{corrige}\end{blocLivre}
\begin{blocLivre}{cor-1-5}\begin{corrige}{5}{Fractions décimales}\demarche{Le dénominateur indique le rang du dernier chiffre.}
$\res{0{,}7}$ ; $\res{0{,}45}$ ; $\res{0{,}003}$ ; $\res{4{,}107}$, car il n’y a aucun centième.
\end{corrige}\end{blocLivre}
\begin{blocLivre}{cor-1-6}\begin{corrige}{6}{Valeur des chiffres}\demarche{Associer chaque chiffre à son rang.}
5 dizaines : 50 ; 8 unités : 8 ; 0 dixième ; 3 centièmes : $0{,}03$ ; 6 millièmes : $0{,}006$. Contrôle : $50+8+0{,}03+0{,}006=\res{58{,}036}$.
\end{corrige}\end{blocLivre}
\begin{blocLivre}{cor-1-7}\begin{corrige}{7}{Le même nombre ?}\demarche{Écrire chaque proposition avec une virgule.}
Conviennent : $\res{\frac{34}{10}}$, $\res{\frac{340}{100}}$ et $\res{3+\frac4{10}}$. En revanche $3+\frac4{100}=3{,}04$ et $\frac{34}{100}=0{,}34$.
\end{corrige}\end{blocLivre}
\begin{blocLivre}{cor-1-8}\begin{corrige}{8}{Nombres mixtes}\demarche{Extraire les unités entières du numérateur.}
$\frac{2\,375}{1\,000}=\res{2+\frac{375}{1\,000}}=2{,}375$ ; $\frac{1\,204}{100}=\res{12+\frac4{100}}=12{,}04$ ; $\frac{89}{10}=\res{8+\frac9{10}}=8{,}9$.
\end{corrige}\end{blocLivre}
\begin{blocLivre}{cor-1-9}\begin{corrige}{9}{Décimal ou pas ?}\demarche{Multiplier numérateur et dénominateur pour obtenir 10, 100 ou 1 000.}
$\frac34=\res{\frac{75}{100}}=0{,}75$ ; $7=\res{\frac{70}{10}}$ ; $\frac15=\res{\frac2{10}}=0{,}2$ ; $\frac2{25}=\res{\frac8{100}}=0{,}08$.
\end{corrige}\end{blocLivre}
\begin{blocLivre}{cor-1-10}\begin{corrige}{10}{Distances dans l’espace}\demarche{Compter en millions pour comparer.}
$150$ millions : $\res{150\,000\,000}$ km. $4{,}5$ milliards, c’est 4 milliards et 500 millions : $\res{4\,500\,000\,000}$ km. En millions : $4\,500:150=\res{30}$ fois.
\end{corrige}\end{blocLivre}
\begin{blocLivre}{cor-1-11}\begin{corrige}{11}{Huit milliards d’humains}\demarche{Un milliard vaut mille millions.}
$\res{8\,000\,000\,000}$ et $\res{4\,000\,000\,000}$. L’écart vaut 4 milliards, soit $\res{4\,000}$ millions de personnes. La population a été multipliée par $\res2$ en 48 ans.
\end{corrige}\end{blocLivre}
\begin{blocLivre}{cor-1-12}\begin{corrige}{12}{Des boîtes de mille}\demarche{Lire un nombre de milliers, puis de millions.}
Le nombre de milliers de vis donne les boîtes : $\res{2\,450\,000}$ boîtes. Un carton contient un million de vis : le nombre de millions donne $\res{2\,450}$ cartons.
\end{corrige}\end{blocLivre}
\begin{blocLivre}{cor-1-13}\begin{corrige}{13}{Le nombre mystère}\demarche{Remplir les rangs un à un.}
Unités 3, dixièmes 4, centièmes 0, millièmes $2\times4=8$ : le nombre est $\res{3{,}408}$, soit $\res{\frac{3\,408}{1\,000}}$.
\end{corrige}\end{blocLivre}
\begin{blocLivre}{cor-1-14}\begin{corrige}{14}{La notation de Stevin}\demarche{Associer chaque indice à un rang.}
$5{,}306$ s’écrit 5\stevin{0}3\stevin{1}0\stevin{2}6\stevin{3}. Le nombre 3\stevin{0}0\stevin{1}9\stevin{2} a 3 unités, 0 dixième et 9 centièmes : $\res{3{,}09}$.
\end{corrige}\end{blocLivre}
\begin{blocLivre}{cor-1-15}\begin{corrige}{15}{Vrai ou faux}\demarche{Distinguer le chiffre et le nombre de dixièmes.}
Elle a $\res{\text{tort}}$. $6{,}25$ contient $\res{62}$ dixièmes (et encore 5 centièmes), et en tout $\res{625}$ centièmes. Le chiffre des dixièmes est 2, mais le nombre de dixièmes est 62.
\end{corrige}\end{blocLivre}
